\documentclass[11pt]{article}
\usepackage[utf8]{inputenc}
\usepackage{amsmath}
\usepackage{geometry}
\usepackage{listings}
\usepackage{xcolor}

\geometry{letterpaper, margin=1in}

\definecolor{codegreen}{rgb}{0,0.6,0}
\definecolor{backcolour}{rgb}{0.95,0.95,0.92}
\definecolor{codepurple}{rgb}{0.58,0,0.82}
\definecolor{codegray}{rgb}{0.5,0.5,0.5}

\lstdefinestyle{mystyle}{
    backgroundcolor=\color{backcolour},   
    commentstyle=\color{codegreen},
    keywordstyle=\color{blue},
    numberstyle=\tiny\color{codegray},
    stringstyle=\color{codepurple},
    basicstyle=\ttfamily\footnotesize,
    breaklines=true,                 
    numbers=left,                    
    numbersep=5pt,                  
    tabsize=2,
    showstringspaces=false
}

\lstset{style=mystyle}

\title{\textbf{Editorial: Problem E - Jayden's Automaton Factory}}
\author{Setter: Hudson}
\date{}

\begin{document}

\maketitle

\section*{Problem Overview}
We are given a board width $W \le 6$, a DFA describing valid piece sequences, and a target length $L \le 10^{18}$. We need to count paths in the DFA such that the Tetris board ends up empty.

\section*{Key Challenges}
\begin{enumerate}
    \item \textbf{Infinite Height vs Finite State:} While the stack can grow infinitely high in theory, we only care if the board \textit{eventually} clears. This suggests we should track the \textbf{relative profile} (skyline) of the board rather than absolute heights.
    \item \textbf{Large $L$:} Since $L$ is up to $10^{18}$, a simulation $O(L)$ is impossible. We must use Matrix Exponentiation $O(\log L)$.
    \item \textbf{Complex State:} The state must combine the Automaton Node and the Tetris Board Profile.
\end{enumerate}

\section*{Solution Approach}

\subsection*{1. Defining the Profile}
Since $W$ is very small ($W \le 6$), we can represent the surface of the current stack as a tuple of heights $(h_1, h_2, \dots, h_W)$.
Crucially, absolute height doesn't affect placement logic—only relative height does.
We normalize the profile by subtracting the minimum height from all columns:
$$ P = (h_1 - \min(h), h_2 - \min(h), \dots, h_W - \min(h)) $$
The "Target State" is the flat profile $(0, 0, \dots, 0)$.

\textbf{Bound on Profile Height:} 
Does the difference between max and min height grow indefinitely?
If we are only searching for paths that return to the flat state $(0,\dots,0)$, we can prune any state where $\max(P) > K$.
For $W=6$, a heuristic limit of $K=8$ is sufficient. If the "hole" becomes deeper than the longest piece (plus a buffer), it becomes impossible to fill that hole without creating new holes higher up.

\subsection*{2. Precomputing Transitions (The "Super-Node")}
Our total system state is defined by:
$$ S = (\text{DFA Node } u, \text{Profile } P) $$
The number of valid profiles for $W=6$ is surprisingly small (a few hundred). The number of DFA nodes is $\le 50$.
Total states $\approx 50 \times (\text{small constant})$.

We perform a BFS/DFS to discover reachable states starting from $(\text{Node } 1, \text{Profile } \{0\dots0\})$.
For each reachable state $(u, P)$ and each outgoing edge in the DFA $u \to v$ with piece $(T, C)$:
\begin{enumerate}
    \item Simulate dropping piece $T$ at column $C$ onto profile $P$.
    \item Calculate new column heights.
    \item If a row is full (all $W$ cols increased), simulate line clear (subtract 1 from all).
    \item Re-normalize (subtract min height).
    \item If max height exceeds cutoff, discard.
    \item Otherwise, we have a valid transition to $(v, P_{new})$.
\end{enumerate}

\subsection*{3. Matrix Exponentiation}
Let $\mathcal{S}$ be the set of valid composite states. Map each state to an index $0 \dots M-1$.
Construct an adjacency matrix $\mathbf{A}$ of size $M \times M$, where $\mathbf{A}[i][j] = 1$ if there is a transition from state $i$ to state $j$.
Note: There may be multiple edges between DFA nodes; sum them up in the matrix.

The answer is the entry corresponding to $(\text{Start Node}, \text{Flat}) \to (\text{Any Node}, \text{Flat})$ in the matrix $\mathbf{A}^L$.
Strictly speaking, we want the count of paths ending at specific states.
We initialize a row vector $V_0$ with $1$ at index $(\text{Node } 1, \text{Flat})$ and $0$ elsewhere.
We compute $V_L = V_0 \times \mathbf{A}^L$.
The final answer is the sum of $V_L[k]$ for all $k$ where the profile part of state $k$ is Flat.

\section*{Complexity}
\begin{itemize}
    \item \textbf{Profile Generation:} The number of profiles is related to the "roughness" of the surface. For $W=6$, this is manageable.
    \item \textbf{Matrix Size:} Let $S_{total} \approx 2000$. Matrix multiplication is $O(S_{total}^3)$. With $L=10^{18}$, we need $\approx 60$ multiplications.
    \item \textbf{Optimization:} $2000^3$ is too slow for 5.0s. However, the matrix is sparse.
    \item \textbf{Actual Constraints:} With $W \le 6$, the number of relevant "Flat-reachable" profiles is usually $< 100$. With $N=50$, matrix size is roughly 5000. This is still tight. We can optimize by noting that for specific $W$, many profiles are isomorphic or unreachable.
\end{itemize}

\newpage
\section*{Implementation (C++)}

\begin{lstlisting}[language=C++]
#include <iostream>
#include <vector>
#include <string>
#include <map>
#include <queue>
#include <algorithm>

using namespace std;

const int MOD = 1e9 + 7;

// --- TETRIS LOGIC ---

struct PieceShape {
    int h, w;
    vector<string> grid;
};

// Definitions matching the problem description
map<char, PieceShape> SHAPES = {
    {'I', {1, 4, {"####"}}},
    {'O', {2, 2, {"##", "##"}}},
    {'T', {2, 3, {".#.", "###"}}},
    {'L', {2, 3, {"..#", "###"}}},
    {'J', {2, 3, {"#..", "###"}}},
    {'S', {2, 3, {".##", "##."}}},
    {'Z', {2, 3, {"##.", ".##"}}}
};

// Represents the "Surface" of the board.
using Board = vector<string>; 

// Max height allowed before we declare state "dead"
const int MAX_PRUNE_HEIGHT = 8; 

// Simulate dropping a piece. Returns empty board if invalid/overflow.
Board simulate(Board b, char pType, int col, int W) {
    PieceShape p = SHAPES[pType];
    int c_idx = col - 1; 

    // Pad board with empty space to simulate drop
    int current_h = b.size();
    int start_y = current_h + 3; // Start safely above
    
    // Find lowest valid position
    int final_y = -1;
    
    for (int y = start_y; y >= 0; --y) {
        bool collision = false;
        // Check collision at this position
        for (int pr = 0; pr < p.h; ++pr) {
            for (int pc = 0; pc < p.w; ++pc) {
                if (p.grid[pr][pc] == '#') {
                    int abs_r = y + (p.h - 1 - pr); 
                    if (abs_r < 0) { collision = true; break; }
                    if (abs_r < current_h && b[abs_r][c_idx + pc] == '#') { 
                        collision = true; break; 
                    }
                }
            }
            if(collision) break;
        }
        
        if (collision) {
            final_y = y + 1;
            break;
        }
        if (y == 0) final_y = 0; // Reached floor
    }

    // Place piece
    for (int pr = 0; pr < p.h; ++pr) {
        for (int pc = 0; pc < p.w; ++pc) {
            if (p.grid[pr][pc] == '#') {
                int abs_r = final_y + (p.h - 1 - pr);
                int abs_c = c_idx + pc;
                while (b.size() <= abs_r) b.push_back(string(W, '.'));
                b[abs_r][abs_c] = '#';
            }
        }
    }

    // Clear lines
    vector<string> next_b;
    for (const string& row : b) {
        bool full = true;
        for (char c : row) if (c == '.') full = false;
        if (!full) next_b.push_back(row);
    }
    
    // Prune logic
    if (next_b.size() > MAX_PRUNE_HEIGHT) return {"INVALID"};
    
    return next_b;
}

// --- MATRIX LOGIC ---

typedef vector<vector<long long>> Matrix;

Matrix multiply(const Matrix& A, const Matrix& B) {
    int n = A.size();
    Matrix C(n, vector<long long>(n, 0));
    for (int i = 0; i < n; ++i)
        for (int k = 0; k < n; ++k)
            if (A[i][k])
                for (int j = 0; j < n; ++j)
                    C[i][j] = (C[i][j] + A[i][k] * B[k][j]) % MOD;
    return C;
}

Matrix power(Matrix A, long long p) {
    int n = A.size();
    Matrix Res(n, vector<long long>(n, 0));
    for (int i = 0; i < n; ++i) Res[i][i] = 1;
    while (p > 0) {
        if (p & 1) Res = multiply(Res, A);
        A = multiply(A, A);
        p >>= 1;
    }
    return Res;
}

// --- MAIN SOLVER ---

struct Edge {
    int v;
    char type;
    int col;
};

int main() {
    ios_base::sync_with_stdio(false);
    cin.tie(NULL);

    int W, N;
    long long L;
    if (!(cin >> W >> N >> L)) return 0;

    int M;
    cin >> M;
    vector<vector<Edge>> adj(N + 1);
    for (int i = 0; i < M; ++i) {
        int u, v, c;
        char t;
        cin >> u >> v >> t >> c;
        adj[u].push_back({v, t, c});
    }

    // BFS to find reachable states (Node, Board)
    map<pair<int, Board>, int> state_to_id;
    vector<pair<int, Board>> id_to_state;
    queue<pair<int, Board>> q;

    Board empty_board; // size 0
    state_to_id[{1, empty_board}] = 0;
    id_to_state.push_back({1, empty_board});
    q.push({1, empty_board});

    int state_cnt = 0;
    
    while(!q.empty()) {
        auto [u, board] = q.front();
        q.pop();
        
        // Try all edges
        for (auto& e : adj[u]) {
            Board next_b = simulate(board, e.type, e.col, W);
            if (next_b.size() == 1 && next_b[0] == "INVALID") continue;
            
            pair<int, Board> next_state = {e.v, next_b};
            if (state_to_id.find(next_state) == state_to_id.end()) {
                state_to_id[next_state] = ++state_cnt;
                id_to_state.push_back(next_state);
                q.push(next_state);
            }
        }
    }

    int total_states = state_cnt + 1;
    Matrix Mat(total_states, vector<long long>(total_states, 0));

    // Fill Matrix
    for (int i = 0; i < total_states; ++i) {
        auto [u, board] = id_to_state[i];
        for (auto& e : adj[u]) {
            Board next_b = simulate(board, e.type, e.col, W);
            if (next_b.size() == 1 && next_b[0] == "INVALID") continue;
            
            int j = state_to_id[{e.v, next_b}];
            Mat[i][j]++;
        }
    }

    // Exponentiate
    Matrix FinalMat = power(Mat, L);

    long long ans = 0;
    for (int i = 0; i < total_states; ++i) {
        // We want transitions 0 -> i where i is a Flat board
        if (FinalMat[0][i] > 0) {
            if (id_to_state[i].second.empty()) {
                ans = (ans + FinalMat[0][i]) % MOD;
            }
        }
    }

    cout << ans << endl;
    return 0;
}
\end{lstlisting}

\end{document}
